\subsection{Probing Large Extra Dimensions with Neutrinos --- arXiv:hep-ph/9904211}

\textbf{Authors:} Gia Dvali, Alexei Yu. Smirnov

\paragraph{主な主張}
brane に局在した左巻きニュートリノと bulk singlet fermion の KK tower の mixing による vacuum oscillation と matter resonance conversion を解析し、多重 KK resonance による solar $\nu_e$ の変換で solar neutrino problem を説明し得ること、さらに少なくとも一つの余剰次元半径が $0.06$--$0.1\,\mathrm{mm}$ 程度になることを示した \cite{Dvali:1999LEDnu}。

\paragraph{新規性}
large extra dimension による体積抑制された neutrino mass の機構を、無限 KK tower との oscillation と varying-density medium での多重 resonance conversion という具体的な neutrino phenomenology へ接続した点が新しい。

\paragraph{位置付け}
LED neutrino phenomenology の基礎論文の一つであり、higher KK modes による複数の matter resonance が現れること自体はすでにこの段階で議論されているため、現在の $T^2$--IceCube 研究では resonance の存在そのものを新規性とはできない。

\paragraph{コメント}
我々の研究で新しい点は、generic $T^2$ では $M_{m,n}^2\propto |m-n\tau|^2/\tau_2$ であるため $E_{\rm res}^{(m,n)}\propto M_{m,n}^2$ として resonance の位置・間隔・縮退構造が complex structure modulus $\tau$ に依存し、その compactification geometry の情報が Earth-matter neutrino propagation に現れることを調べる点にある。

UV の観点では、本論文は二つの大きな余剰次元について $\Delta\sim \xi^2\sum_{n,k}[n^2+(R/R')^2k^2]^{-1}$ を考え、過程の energy release $Q$ まで生成可能な KK states を積分すると $\Delta_Q\propto \xi^2(R'/R)\ln(QR')$ となることを示しているので、二次元 KK tower の logarithmic enhancement 自体は既知である。一方、ここでの $Q$ は主として kinematic production cutoff であり、generic $T^2$ の 6D EFT における UV cutoff、brane-localized operators、current normalization を含む UV matching / renormalization は扱われていないため、我々の計算で現れる absolute active weight と full probability の cutoff dependence を定量予言へ昇格させるには別途 UV matching が必要である。
